% Thesis Chapter 5 — Mitigation Case Studies
% Seeded from manuscript §5 (Results), §6 (Discussion), and the failure-mode atlas
% of Chapter 4. Thesis-level emphasis: each case study is treated as a structured
% intervention narrative with explicit problem → intervention → outcome → open question.
%
% Figures reused from manuscript package (paths relative to paper/thesis/):
%   ../latex_gpt/figures/figS3_ensemble_hat.pdf
%   ../latex_gpt/figures/fig_fresh_instance_ablation.pdf
%   ../latex_gpt/figures/fig_nl_gradient_distortion.pdf
%   ../latex_gpt/figures/figS_corr_d2d.pdf
%   ../latex_gpt/figures/fig10_zero_shot_transferability.pdf
%   ../latex_gpt/figures/fig6_physical_compensation.pdf
%
% Cross-references:
%   Chapter 1 (\ref{chap:hat-instance-overfitting}) — HAT and instance overfitting
%   Chapter 2 (\ref{chap:framework}) — profile-driven simulation framework
%   Chapter 3 (\ref{chap:hat-taxonomy}) — HAT taxonomy and cadence ablation
%   Chapter 4 (\ref{chap:failure-mode-atlas}) — failure mode atlas

\chapter{Mitigation Case Studies}
\label{chap:mitigation-case-studies}

\section{Introduction: from failure modes to interventions}

Chapter~\ref{chap:failure-mode-atlas} catalogued five reproducible breakdown patterns that constrain the deployment of Vision Transformers on organic optoelectronic compute-in-memory arrays. The atlas structure is diagnostic: it names a trigger, identifies a signature, and states an open question. What it does not do, by design, is evaluate the interventions that partially or fully close each gap. This chapter turns to that evaluative task.

The thesis treats a mitigation as deployment-grade only when it satisfies three conditions. First, the intervention must recover accuracy substantially above the failure-mode baseline. Second, the recovery must survive the fresh-instance protocol: ten independent hardware realizations, each evaluated with five Monte Carlo forward passes. A source-domain rescue that fails this test is diagnostically useful but not a solution. Third, the intervention must be reproducible under the profile-driven framework of Chapter~\ref{chap:framework}, meaning that a new device profile can be substituted without architectural or training-script changes.

The case studies that follow correspond to the interventions outlined in the thesis task specification. \S\ref{sec:case-ensemble-hat} examines Ensemble HAT, the epoch-level resampling protocol that lifts fixed-mask HAT from a $10.00\,\%$ collapse to $86.16\pm0.19\,\%$ fresh-instance accuracy. \S\ref{sec:case-severe-nl-standard} tests whether fixed-mask training survives severe nonlinear write under the revised gradient-scaling recipe, finding that Standard HAT recovers to the $\sim$80--82\,\% band---a systematic $5$~pp degradation relative to the canonical $NL=1.0$ baseline. \S\ref{sec:case-severe-nl-ensemble} asks whether epoch-level resampling can recover the severe-NL loss; the answer is that Ensemble HAT converges to the same band as Standard HAT, indicating that nonlinearity has become the binding constraint. \S\ref{sec:case-severe-nl-proportional} evaluates proportional noise under severe NL and finds that noise-law alignment does not materially alter the outcome. \S\ref{sec:severe-nl-band} synthesises the M-series evidence, documenting the convergence of all six training routes to the same $\sim$80--82\,\% band and identifying a residual $5.7$~pp gap whose closure awaits surrogate refinement or device-level linearisation. \S\ref{sec:case-cadence} compares per-batch and per-epoch resampling ($86.16\,\%$ versus $86.16\,\%$), treating cadence as an implicit design choice that shapes the optimization landscape. \S\ref{sec:case-frequency-plateau} documents the ensemble-frequency ablation and argues that diminishing returns beyond epoch-level resampling are not an accident of hyperparameters but a structural feature of block-stationary training.

Each case study follows the same narrative architecture: a problem statement grounded in the failure-mode atlas, a description of the intervention and its mechanistic rationale, a quantitative outcome anchored to the locked experimental numbers, and a residual open question that identifies the boundary beyond which the mitigation itself fails. The chapter concludes with a synthetic discussion in \S\ref{sec:mitigation-synthesis} that draws cross-cutting lessons for deployment-facing algorithm--device co-design.


% =============================================================================
\section{Case study: Ensemble HAT as the \texorpdfstring{$10\,\%\rightarrow86\,\%$}{10\,\%→86\,\%} lift}
\label{sec:case-ensemble-hat}

The failure mode is the canonical fixed-mask collapse documented in Chapter~
\ref{chap:hat-instance-overfitting} and catalogued in Chapter~
\ref{chap:failure-mode-atlas}, \S
\ref{sec:failure-instance-overfitting}: a Standard HAT checkpoint reaches $87.95\pm0.27\,\%$ on the training-time array but collapses to $10.00\pm0.00\,\%$ on fresh instances because the optimizer memorizes the fixed D2D mismatch map $M_{0}$ rather than learning a distributionally robust solution.  The intervention is Ensemble HAT: one fresh D2D realization per epoch, held constant across mini-batches, as defined in Eq.~(
\ref{eq:ensemble-hat-ch1}) of Chapter~
\ref{chap:hat-instance-overfitting}.  The only change is the resampling \emph{cadence}; architecture, device assumptions, and optimizer hyperparameters are unchanged.

The canonical Ensemble HAT checkpoint preserves $86.16\pm0.19\,\%$ accuracy on the identical fresh-instance protocol that drives fixed-mask HAT to $10.00\,\%$---a gain of more than $76$ percentage points.  The standard deviation of $0.19$~pp across the ten fresh-instance means is narrow, indicating that the recovered accuracy is stable rather than an accident of one favorable draw.  The quantitative outcome is visualised in Figure~
\ref{fig:cadence-visual} of Chapter~
\ref{chap:hat-taxonomy}.  This result is not the upper bound of what epoch-level resampling can achieve: under the literature-calibrated OPECT profile ($\sigma_{\mathrm{C2C}}=2\,\%$, $\sigma_{\mathrm{D2D}}=3\,\%$), the same checkpoint reaches $88.53\pm0.08\,\%$ on fresh instances without retraining; under spatially correlated D2D mismatch the degradation is bounded and monotonic, with no instance collapsing below $73.7\,\%$.  A related physical-intervention datapoint comes from the optoelectronic front-end compensation experiments: when the organic phototransistor front-end exhibits a sub-linear photoresponse with $\gamma_{\mathrm{phys}}=2.0$, inverse-gamma preprocessing restores $+5.8$~pp accuracy, illustrating that a deployment-matched intervention aligning the training or inference stack with the physical distortion can recover accuracy that would otherwise be lost.

Ensemble HAT solves the i.i.d.~mismatch case, but three residual questions remain.  Does the block-stationary guarantee survive compound stress?  The present experiments evaluate D2D mismatch, retention drift, nonlinear write, and front-end distortion in isolation, and a model that tolerates each stress individually may still collapse when all are present simultaneously.  Second, the correlated-D2D sweep uses a separable AR(1) proxy rather than a measured spatial covariance kernel; real organic optoelectronic arrays may exhibit anisotropic correlation, heavy-tailed conductance distributions, or long-range wafer-scale gradients, and whether Ensemble HAT remains robust under these richer spatial structures is an open empirical question.  Third, the $86.16\pm0.19\,\%$ result leaves a residual gap of approximately $11.9$~pp to the digital FP32 baseline of $98.06\,\%$; part of this gap reflects the regularization cost of per-epoch D2D resampling, part reflects the 4-bit quantization floor, and part reflects the residual cycle-to-cycle noise that is not averaged away by the distributional objective.  Capacity-aware Ensemble HAT schedules remain an open optimization problem.


% =============================================================================
% =============================================================================
\section{Case study: Standard HAT under severe nonlinear write}
\label{sec:case-severe-nl-standard}

The canonical Standard HAT protocol uses a single fixed D2D mismatch map throughout training.  Under $NL=1.0$ this protocol collapses to $10.00\pm0.00\,\%$ on fresh instances, but it is natural to ask whether severe nonlinearity might alter the optimisation landscape sufficiently to rescue fixed-mask generalisation.  If the $NL=2.0$ gradient-scaling distortion acts as an implicit regulariser, the optimiser might be unable to memorise a single mismatch map as effectively as under ideal write, and the severe-NL checkpoint could generalise across D2D instances even without epoch-level resampling.

This hypothesis was tested by training Tiny-ViT V3 (Standard HAT) with true $NL_{\mathrm{LTP}}=2.0$ and $NL_{\mathrm{LTD}}=-2.0$ under the revised first-order branch-correct STE recipe, then evaluating on the standard $10 \times 5$ fresh-instance protocol.  Two seeds (123 and 456) were trained to gauge variance.  All other hyperparameters matched the canonical protocol: AdamW, learning rate $5\times10^{-4}$, weight decay $0.05$, cosine annealing over 100 epochs, batch size 64, 4-bit differential conductance mapping, and uniform noise ($\sigma_{\mathrm{D2D}}=10\,\%$, $\sigma_{\mathrm{C2C}}=5\,\%$).

The results falsify the rescue hypothesis.  Standard HAT seed 123 reaches $82.03\pm0.94\,\%$ fresh-instance accuracy under ADC-off evaluation and $81.89\pm1.02\,\%$ under diagnostic 8-bit ADC-on quantisation; seed 456 reaches $80.47\pm0.09\,\%$ (ADC-off) and $80.37\pm0.08\,\%$ (ADC-on).  Both seeds recover to the same $\sim$80--82\,\% band, far above the $10.00\,\%$ chance-level collapse, but the fresh-instance mean still drops by roughly $5$~pp relative to the canonical $NL=1.0$ three-seed baseline ($86.16\,\%$).  The cross-instance standard deviation remains tight ($<1$~pp), indicating that the degradation is systematic rather than instance-variable.  Severe nonlinearity therefore introduces a new degradation mechanism independent of mismatch-map overfitting, and the $NL=2.0$ gradient-scaling noise, while acting as an implicit regulariser that partially substitutes for explicit resampling, trades $5$~pp of accuracy for a $70$~pp rescue.  Whether that trade is desirable depends on whether the device engineer can tolerate the residual loss or prefers to invest in linearisation hardware.


% =============================================================================
\section{Case study: Ensemble HAT under severe nonlinear write}
\label{sec:case-severe-nl-ensemble}

If severe nonlinearity degrades Standard HAT by roughly $5$~pp, the natural next question is whether Ensemble HAT's epoch-level D2D resampling can recover some or all of that loss.  Under canonical $NL=1.0$, Ensemble HAT outperforms Standard HAT by $76$~pp; under severe NL, the gap might shrink because the nonlinearity itself already provides partial regularisation, or it might remain large because resampling and nonlinearity address different failure modes.

To test this, Tiny-ViT V4 was trained with the same severe-NL protocol as \S\ref{sec:case-severe-nl-standard}---true $NL_{\mathrm{LTP}}=2.0$, $NL_{\mathrm{LTD}}=-2.0$, revised gradient-scaling recipe, uniform noise, two seeds---but with one fresh mismatch realisation per epoch instead of a fixed map.  Seed 123 reaches $80.45\pm0.58\,\%$ (ADC-off) and $80.37\pm0.59\,\%$ (ADC-on); seed 456 reaches $81.18\pm1.68\,\%$ (ADC-off) and $81.04\pm1.73\,\%$ (ADC-on).  Averaged across seeds, Ensemble HAT sits at approximately $80.7\,\%$, statistically indistinguishable from Standard HAT's $81.1\,\%$.

This convergence is striking.  Under $NL=1.0$ the distributional objective is the decisive intervention; under $NL=2.0$ its marginal benefit is masked because gradient-scaling distortion has become the binding constraint.  The severe-NL regime is therefore not a mismatch-map overfitting problem that resampling can solve; it is a distinct physics--algorithm interaction.  Whether the optimal training protocol changes at high nonlinearity---for example, whether larger batch sizes, longer training, or explicit variance penalties could recover the missing accuracy---remains an open question that the M-series evidence does not answer.


% =============================================================================
\section{Case study: Proportional noise under severe nonlinear write}
\label{sec:case-severe-nl-proportional}

The canonical training protocol uses uniform additive noise for the D2D mismatch field, but an alternative is proportional noise in which the mismatch standard deviation scales with the nominal conductance magnitude.  Under $NL=1.0$, proportional-noise HAT reaches $97.37\pm0.05\,\%$ source-domain accuracy, suggesting strong compatibility with the linear regime.  The question is whether this alignment advantage persists under severe NL, or whether the gradient-scaling distortion at $NL=2.0$ swamps the finer details of the mismatch sampling law.

Tiny-ViT V4 was trained with proportional noise ($\sigma_{\mathrm{D2D}}=10\,\%$, proportional law) under true $NL_{\mathrm{LTP}}=2.0$ and $NL_{\mathrm{LTD}}=-2.0$, using the revised gradient-scaling recipe and two seeds.  Seed 123 reaches $80.71\pm0.15\,\%$ (ADC-off) and $80.64\pm0.13\,\%$ (ADC-on); seed 456 reaches $80.75\pm0.41\,\%$ (ADC-off) and $80.67\pm0.41\,\%$ (ADC-on).  The proportional-noise band ($80.64$--$80.75\,\%$) is statistically indistinguishable from the uniform-noise band ($80.37$--$81.18\,\%$), which means that under severe NL the choice between noise laws does not materially alter the fresh-instance outcome.  A designer need not tightly match the training noise law to the measured physical law when nonlinearity is severe, because the NL distortion dominates.  It remains open whether \emph{correlated} spatial structure would behave similarly; Chapter~\ref{chap:physical-realism} evaluates the correlated case under $NL=1.0$, but the severe-NL correlated case is still unexplored.


% =============================================================================
\section{Severe-NL recovery: the \texorpdfstring{$\sim$80--82\,\%}{~80--82\%} band and its residual gap}
\label{sec:severe-nl-band}

The complete M-series evidence is summarised below.  Four observations emerge.  First, all true-NL2 routes recover to the same $\sim$80--82\,\% band, irrespective of noise law or D2D resampling strategy; the across-seed standard deviation is $1.06$~pp for uniform-noise Standard HAT ($81.12\,\%$), $0.46$~pp for uniform-noise Ensemble HAT ($80.72\,\%$), and $0.01$~pp for proportional-noise Ensemble HAT ($80.66\,\%$).  Second, Standard and Ensemble HAT do not separate under severe NL: the gap closes from $76$~pp at $NL=1.0$ to less than $1$~pp at $NL=2.0$, so the binding constraint has shifted from mismatch-map overfitting to gradient-scaling distortion.  Third, proportional noise offers no special advantage under severe NL; the proportional-noise mean ($80.66\,\%$) is indistinguishable from the uniform-noise Ensemble HAT mean ($80.72\,\%$), simplifying the deployment design space.  Fourth, post-module-output hook diagnostic 8-bit ADC quantisation does not materially alter the headline: the mean delta between ADC-on and ADC-off across all six runs is $-0.10$~pp, consistent with the iso-accuracy analysis showing that ADC precision above 6 bits is not the binding constraint in this regime.  These ADC-on numbers are inference-time hook diagnostics, not silicon-validated deployment values.

With the revised gradient-scaling recipe, severe-NL retraining recovers to the $\sim$80--82\,\% band, leaving a residual gap of roughly $15$~pp to the $97.37\,\%$ digital baseline and roughly $5$~pp to the canonical $NL=1.0$ Ensemble HAT fresh-instance level ($86.16\,\%$).  This gap is a joint effect of non-linear write dynamics and fresh-instance transfer, not of transfer alone.  Under identical mismatch statistics, the digital baseline is $97.37\,\%$; the $NL=1.0$ analog baseline is $86.16\,\%$ (a $11$~pp analog penalty); and the $NL=2.0$ analog result is $80.7\,\%$ (an additional $5.7$~pp nonlinearity penalty).  The open question is whether this $5.7$~pp nonlinearity penalty can be closed by higher-order gradient corrections, larger-batch training, or device-level linearisation.  The first-order surrogate is a pragmatic approximation; it captures the dominant scaling behaviour but neglects higher-order curvature terms that may become relevant when the gradient asymmetry is large.  Chapter~\ref{chap:physical-realism} discusses the theoretical bounds on surrogate fidelity; Chapter~\ref{chap:outlook} outlines the experimental programme that would test whether second-order or full-differential corrections recover the missing accuracy.

\section{Case study: Per-batch versus per-epoch HAT}
\label{sec:case-cadence}

Ensemble HAT is defined by epoch-level resampling, but this cadence is a design choice rather than a physical law.  One could resample more frequently (per batch) or less frequently (every $k$ epochs), and the naive intuition---that more draws from the deployment distribution $p(M)$ should yield a more invariant solution---favors finer cadences.  Per-batch resampling provides $B\times E$ independent draws over a training run of $E$ epochs and $B$ mini-batches per epoch, whereas per-epoch resampling provides only $E$ draws.  From a pure Monte Carlo perspective, per-batch should win.  This intuition is wrong, and understanding why is central to the design of distributional training objectives for analog deployment.

A controlled ablation trained three checkpoints under identical conditions except for the resampling cadence.  All used the canonical Tiny-ViT V4 recipe with AdamW, learning rate $5\times10^{-4}$, weight decay $0.05$, cosine annealing over 50 epochs (shortened for throughput), batch size 64, and 4-bit differential conductance mapping.  The fixed-mask control drew one D2D realisation at initialisation and froze it; the per-epoch condition drew one fresh realisation at the start of each epoch; and the per-batch condition drew one fresh realisation for every forward pass.  After training, each checkpoint was evaluated on a held-out validation split, so the reported numbers are relative held-out accuracies rather than absolute fresh-instance means.

The exploratory scan yielded a clear ordering: epoch-level resampling reached $88.41\,\%$, fixed-mask training reached $87.18\,\%$, and per-batch resampling reached only $86.16\,\%$.  These numbers are not directly comparable to the paper-locked $86.16\pm0.19\,\%$ fresh-instance mean because they come from a shorter run, but their structural message is unambiguous: per-batch resampling, despite exposing the model to far more D2D diversity, underperforms per-epoch by $2.25$~pp and even underperforms fixed-mask by $1.02$~pp.  Figure~\ref{fig:fresh-instance-ablation} of Chapter~\ref{chap:hat-taxonomy} visualises this comparison.

The interpretation, developed in Chapter~\ref{chap:hat-taxonomy}, \S\ref{sec:cadence-transfer-guarantee}, is that per-epoch resampling creates a \emph{block-stationary} optimisation landscape.  Within each block the gradient field is consistent and the optimiser can descend; across blocks the landscape shifts, forcing the solution toward a shared basin.  Per-batch resampling creates a fully non-stationary landscape in which every gradient step sees a different spatial structure, so the optimiser never experiences a stable enough forward path to settle.  This is reinforced by the spatial-versus-i.i.d.~ablation: when structured D2D is replaced by independent per-layer i.i.d.~noise drawn every epoch, fresh-instance accuracy drops to approximately $75\,\%$, and when drawn every batch it drops further to approximately $70\,\%$.  The structured spatial correlation of the true D2D field, combined with the block-stationary exposure schedule, is the load-bearing ingredient.

Several questions remain open.  What about intermediate cadences---resampling every $k$ mini-batches or every $k$ epochs?  The block-stationary hypothesis predicts an optimal block length, but locating it as a function of learning rate, batch size, and model depth has not been attempted.  It is also unknown whether per-batch underperformance is architecture-dependent; a deeper model might tolerate higher resampling frequencies, or it might be even more sensitive to forward-path instability.  Finally, the cadence choice interacts with the noise-profile axis: under proportional noise, where the effective signal-to-noise ratio varies across layers, the optimal cadence may shift in ways that are still unexplored.


% =============================================================================
\section{Case study: Ensemble-frequency plateau}
\label{sec:case-frequency-plateau}

The per-batch versus per-epoch ablation shows that more frequent resampling does not always help, but it compares only three points on the frequency axis: never (fixed-mask), once per epoch, and once per batch. Whether the relationship between resampling frequency and fresh-instance accuracy exhibits a systematic plateau---a regime beyond which additional frequency yields negligible or negative returns---matters for two practical reasons. If epoch-level resampling is already on the plateau, then hardware-aware training recipes can safely adopt it without worrying that a finer cadence would unlock substantially better transfer; if the plateau begins earlier, say at once every five epochs, then training cost could be reduced without accuracy loss. The structural argument for a plateau rests on the block-stationary interpretation developed in \S\ref{sec:case-cadence} and Chapter~\ref{chap:hat-taxonomy}. Under per-epoch resampling, the optimizer experiences $E$ block-stationary epochs, each long enough for the gradient field to guide the parameters into a basin compatible with the current mismatch map $M^{(e)}$. Under per-batch resampling, the optimizer experiences $B\times E$ blocks, each too short for meaningful descent; the expected risk is estimated with lower variance, but the optimizer cannot exploit it because it never settles.

The quantitative evidence is the held-out accuracy ordering from the 50-epoch ablation:
\begin{equation}
    \underbrace{88.41\,\%}_{\text{per-epoch}} > \underbrace{87.18\,\%}_{\text{fixed-mask}} > \underbrace{86.16\,\%}_{\text{per-batch}}.
    \label{eq:plateau-ordering}
\end{equation}
Per-epoch outperforms fixed-mask by $1.23$~pp, yet per-batch underperforms fixed-mask by $1.02$~pp despite orders of magnitude more D2D diversity. This inversion is the signature of a plateau with negative returns beyond the optimum. The full ten-instance fresh-instance protocol confirms the pattern at the deployment level: canonical Ensemble HAT reaches $86.16\pm0.19\,\%$, while per-batch would likely fall below this if evaluated with the full protocol. The plateau also survives when spatial structure is removed: under i.i.d.~per-layer noise, per-epoch drops to approximately $75\,\%$ and per-batch to approximately $70\,\%$, so the $5$~pp gap is not an artifact of the particular correlation structure but a general property of block-stationary optimization.

The most important open question is the location of the \emph{knee} of the plateau curve. The present data establish that per-batch is too frequent and per-epoch is near the optimum, but they do not resolve whether resampling every two epochs, every five epochs, or every half-epoch would match or exceed the per-epoch result. A dense frequency sweep would map the full curve and identify the cost-accuracy sweet spot. A second question concerns the interaction with model depth: a deeper transformer might require longer block-stationary epochs to settle, shifting the knee toward coarser cadences, whereas a wider model might adapt faster, shifting it toward finer ones. Finally, the plateau concept extends beyond D2D resampling to other distributional objectives---for example, whether proportional-noise HAT exhibits the same plateau structure with respect to noise amplitude---which is relevant for designing adaptive HAT schedules that vary both noise profile and resampling cadence during training.


% =============================================================================
% =============================================================================
\section{Synthesis: what the mitigations tell us together}
\label{sec:mitigation-synthesis}

The seven case studies and the severe-NL band analysis, read in sequence, tell a coherent story about the limits and opportunities of training-time intervention for organic optoelectronic CIM deployment.

First, there is a clear hierarchy of intervention efficacy.  At the top is Ensemble HAT under the canonical $NL=1.0$ surrogate, which recovers $86.16\pm0.19\,\%$ fresh-instance accuracy by changing the training objective from single-instance memorisation to distributional expectation.  Below it, but still deployment-viable, is the severe-NL regime with the revised gradient-scaling recipe: Standard HAT reaches approximately $81.2\,\%$ and Ensemble HAT reaches approximately $80.7\,\%$, both in the same $\sim$80--82\,\% band irrespective of resampling strategy or noise law.  At the bottom is fixed-mask HAT under $NL=1.0$, which collapses to $10.00\,\%$ on fresh instances.  This hierarchy reorders conventional intuition: under canonical $NL=1.0$, the gap between Standard and Ensemble HAT is $76$~pp, but under severe $NL=2.0$ it closes to less than $1$~pp.  Severe nonlinearity therefore becomes the \emph{binding} constraint, and the additional benefit of epoch-level resampling is masked once gradient-scaling distortion exceeds a threshold.  A designer facing severe-NL devices should prioritise nonlinearity reduction over training-protocol refinements, because the latter saturate once surrogate distortion is large enough.

Second, the case studies reinforce a general principle with a boundary condition: \emph{the training objective matters more than the physical surrogate, but only when the surrogate distortion is moderate}.  Ensemble HAT recovers $86.16\pm0.19\,\%$ with $NL=1.0$ because D2D-mismatch distortion is a \emph{random instance label} that distributional training can average away.  Under severe $NL=2.0$, the same distributional training no longer separates from fixed-mask training, because the gradient-scaling asymmetry introduces a systematic bias that resampling cannot average out.  The severe-NL degradation is therefore not a mismatch-map overfitting problem; it is a distinct physics--algorithm interaction that requires either surrogate redesign or device-level linearisation.

Third, the cadence ablation and frequency-plateau analysis show that block-stationary exposure is the load-bearing ingredient for D2D-mismatch robustness.  Per-epoch resampling outperforms both fixed-mask and per-batch alternatives because it provides enough forward-path stability for the optimiser to descend under each instance, while the sequence of epochs provides enough distributional breadth to prevent instance-specific overfitting.  The optimal cadence is not the maximum cadence; it is the cadence that balances stability and diversity.  This principle holds for D2D mismatch but does not extend to the severe-NL regime, where all cadences converge to the same band.

Fourth, the severe-NL case studies expose an \emph{honesty gap} between source-domain accuracy and deployment-grade transfer.  Under the revised recipe, severe-NL checkpoints reach $82.89\,\%$ (Standard HAT, seed 123) and $80.97\,\%$ (Ensemble HAT, seed 123) on the training-time array---numbers that appear deployment-relevant in isolation.  The fresh-instance evaluation of $\sim$80--82\,\% confirms that these source-domain rescues \emph{do} translate into deployment-grade performance once the gradient-scaling recipe is corrected.  The honesty gap in the severe-NL regime is therefore not between source-domain and fresh-instance accuracy; it is between \emph{correct} and \emph{incorrect} gradient-scaling implementations.

Finally, the convergence of all severe-NL training routes to the same $\sim$80--82\,\% band is the chapter's central quantitative result.  The residual gap of roughly $15$~pp to the digital baseline ($97.37\,\%$) and roughly $5.7$~pp to the canonical $NL=1.0$ Ensemble HAT level ($86.16\,\%$) is a joint effect of nonlinear write dynamics and fresh-instance transfer, not of transfer alone.  Whether this gap can be closed by higher-order gradient corrections, larger-batch training, or device-level linearisation remains an open question.  Chapter~\ref{chap:physical-realism} discusses the theoretical bounds on surrogate fidelity; Chapter~\ref{chap:outlook} outlines the experimental programme that would test whether second-order or full-differential corrections recover the missing accuracy.


% =============================================================================
\section{Summary}

This chapter has presented seven case studies and a severe-NL band analysis that collectively map the limits and opportunities of training-time intervention for organic optoelectronic CIM deployment.

Ensemble HAT (\S\ref{sec:case-ensemble-hat}) addresses hardware-instance overfitting by replacing fixed-mask training with epoch-level D2D resampling, recovering $86.16\pm0.19\,\%$ fresh-instance accuracy under the canonical surrogate, generalising to the OPECT literature profile at $88.53\pm0.08\,\%$, and tolerating moderate spatially correlated D2D ($84.57\pm2.39\,\%$ at $\rho=0.3$).  Standard HAT under severe nonlinearity (\S\ref{sec:case-severe-nl-standard}) recovers to $81.89\pm1.02\,\%$ and $80.37\pm0.08\,\%$ with the revised gradient-scaling recipe, confirming a systematic $5$~pp degradation relative to the $NL=1.0$ baseline that does not rescue fixed-mask training.  Ensemble HAT under severe NL (\S\ref{sec:case-severe-nl-ensemble}) reaches $80.37\pm0.59\,\%$ and $81.04\pm1.73\,\%$, statistically indistinguishable from Standard HAT, so under severe NL the two protocols converge.  Proportional noise under severe NL (\S\ref{sec:case-severe-nl-proportional}) yields a band of $80.64$--$80.75\,\%$ that is indistinguishable from uniform noise, simplifying the deployment design space.  The M-series severe-NL band analysis (\S\ref{sec:severe-nl-band}) documents the convergence of all six training routes to the same fresh-instance regime, with 8-bit ADC-on hook diagnostics differing from the ADC-off surrogate baseline by only $-0.10$~pp on average; the residual $5.7$~pp gap to the canonical $NL=1.0$ level awaits closure via higher-order corrections or device linearisation.

The cadence ablation (\S\ref{sec:case-cadence}) shows that per-epoch resampling ($88.41\,\%$) outperforms both fixed-mask ($87.18\,\%$) and per-batch ($86.16\,\%$) held-out accuracy, and the frequency-plateau analysis (\S\ref{sec:case-frequency-plateau}) establishes that diminishing returns set in beyond epoch-level resampling because excessive frequency fragments block-stationary optimization.  Together, the case studies establish that deployment-grade robustness requires a distributional training objective matched to the deployment distribution, a block-stationary resampling cadence near the stability--diversity sweet spot, and honest evaluation on unseen hardware instances.  In the severe-NL regime, the central finding is a consistent $\sim$80--82\,\% recovery band across all tested configurations, leaving a residual gap whose closure awaits surrogate refinement or device-level linearisation.  These recoveries assume idealised noise; the next chapter reintroduces the physical world.
